%%==================================================
%% chapter01.tex for BIT Master Thesis
%% modified by yang yating
%% version: 0.1
%% last update: Dec 25th, 2016
%%==================================================
\chapter{基于回归排序协同强化学习的端到端温压智能优化}
\enchapter{Introduction}
\label{chap:intro}
\section{面向scCO₂制药工艺的温压优化问题建模}
\ensection{Research Background and Motivation}

近年来，随着人们生活水平的不断提高，人们越来越注重周围环境对身体健康的影响。作为服装是人们时时刻刻最贴近的环境，尤其是内衣，对人体健康有很大的影响。由于合时刻刻最贴近的环境，尤其是内衣，对人体健康有很大的影响。由于合成纤维的衣着舒适性、手感性，天然纤维的发展又成为人们关注的一大热点。\supercite{Takahashi1996Structure,Xia2002Analysis,Jiang1989,Mao2000Motion,Feng1998}

\subsection{优化问题形式化描述}
\ensubsection{Research Progress Overview in Home and Abroad}

\subsection{基于单步马尔可夫决策过程的建模}
\ensubsection{Research Progress Overview in Home and Abroad}

\section{SHIFT-TPO端到端智能优化框架设计}
\ensection{Research Progress Overview in Home and Abroad}
%\label{sec:***} 可标注label

\subsection{框架总体架构 }
\ensubsection{Research Progress Overview in Home and Abroad}

\subsection{回归-排序协同奖励模型(RRCN)构建}
\ensubsection{Research Progress Overview in Home and Abroad}

\subsection{基于PPO的策略优化与软约束设计}
\ensubsection{Research Progress Overview in Home and Abroad}

%\label{sec:features}

\section{实验设计与评价体系}
\ensection{Organizational Structure of This Paper}

\subsection{数据集构建与预处理}
\ensubsection{Research Progress Overview in Home and Abroad}

\subsection{对比基线方法}
\ensubsection{Research Progress Overview in Home and Abroad}

\subsection{评价指标体系}
\ensubsection{Research Progress Overview in Home and Abroad}

\subsection{实验环境与训练细节}
\ensubsection{Research Progress Overview in Home and Abroad}

\section{实验结果与模型性能分析}
\ensection{Organizational Structure of This Paper}

\subsection{温压寻优性能对比与分析}
\ensubsection{Research Progress Overview in Home and Abroad}

\subsection{奖励模型精度与泛化能力评估}
\ensubsection{Research Progress Overview in Home and Abroad}

\subsection{核心模块消融实验分析}
\ensubsection{Research Progress Overview in Home and Abroad}

\subsection{模型物理可解释性分析}
\ensubsection{Research Progress Overview in Home and Abroad}

\section{本章小结}
\ensection{Organizational Structure of This Paper}
